\subsection{Simple Modules over an Artinian Ring}

Let A be a ring. By abuse of language, we say that a left ideal \(\alpha\) is a minimal left ideal of A if it is a minimal element of the set of nonzero left ideals of A, ordered by inclusion. We define right minimal ideals and two-sided minimal ideals analogously.

Let \(\mathfrak{a}\) be a left ideal of A. Then \(\mathfrak{a}\) is a simple A-module if and only if it is a minimal left ideal of A.

Every nonzero left ideal of a left Artinian (VIII, p.~1, Definition 1) ring contains a minimal left ideal.

PROPOSITION 4. --- Let A be a ring that has a minimal left ideal a. Every faithful simple A-module is isomorphic to the A-module a.

Let M be a faithful simple A-module. Let \(\alpha\) be a nonzero element of a. Since the A-module M is faithful, there exists an element x of M such that \(\alpha x \neq 0\) . The homomorphism \(a \mapsto ax\) from a to M is then nonzero; it is therefore an isomorphism by Proposition 2 of VIII, p.~47.

PROPOSITION 5. --- Let A be a left Artinian ring and M a faithful A-module.

\begin{enumerate}
\def\labelenumi{\alph{enumi})}
\item
  There exist a natural number \(m\) and an isomorphism from \(\mathbf{A}_s\) to a submodule of \(\mathbf{M}^m\) .
\item
  If M admits a Jordan--Hölder series \((\mathrm{M}_{i})_{0\leqslant i\leqslant n}\) , then every simple A-module is isomorphic to one of the quotients \(M_{i}/M_{i+1}\) .
\end{enumerate}

By VIII, p.~2 applied to the A-module \(A_{s}\) and the annihilators of the elements of M, there exist elements \(x_{1},\ldots,x_{m}\) of M such that the annihilator of M is the intersection of the annihilators of the \(x_{i}\) . Since the A-module M is faithful, the A-linear mapping \(a\mapsto(ax_{1},\ldots,ax_{m})\) from \(A_{s}\) to \(M^{m}\) is injective, whence a).

Under the assumption of b), every simple quotient of \(A_{s}\) is isomorphic to a quotient of a Jordan--Hölder series of \(M^{m}\) (I, §4, No.~7, p.~43, Corollary), hence to one of the modules \(M_{i}/M_{i+1}\) (I, §4, No.~7, p.~43, Theorem 6). We conclude that every simple A-module is isomorphic to a quotient of \(A_{s}\) .

Remark. --- Proposition 5 applies, in particular, to the following two cases:

\begin{enumerate}
\def\labelenumi{\alph{enumi})}
\item
  Let A be an algebra over a commutative field K, and let M be a faithful finite-dimensional module over K. Then M is an A-module of finite length, and the countermodule of M is finitely generated. The ring A is left Artinian (VIII, p.~9, Proposition 6). There exists a Jordan--Hölder series \((\mathrm{M}_{i})_{0\leqslant i\leqslant n}\) of the A-module M, and every simple A-module is isomorphic to one of the modules \(M_{i}/M_{i+1}\) for \(0\leqslant i\leqslant n-1\) .
\item
  Let A be a left Artinian ring. The module \(A_{s}\) has finite length (VIII, p.~6, Theorem 1). Since the A-module \(A_{s}\) has finite length, there exists a decreasing sequence \((\mathfrak{a}_{i})_{0\leqslant i\leqslant n}\) of left ideals of A such that \(a_{0}=A\) and \(a_{n}=0\) and that the A-modules \(S_{i}=\mathfrak{a}_{i-1}/\mathfrak{a}_{i}\) are simple for \(1\leqslant i\leqslant n\) . Then every simple A-module is isomorphic to one of the modules \(S_{1},\ldots,S_{n}\) .
\end{enumerate}

